% Caption draft for paper/figs/fig_window.pdf (figure*, full width). Numbers from research note E9.
\caption{%
Region of the ALP mass $\mu$ and decay constant $f$ in which self-interactions could reconcile a
spin pile-up of first-generation (1G) black holes with the rapid spins of the primary of GW241011
and of GRS~1915+105.
Gray: excluded, because the primary of GW241011 (blue) or GRS~1915+105 (black) would lose more spin
than observed.
Dashed: above each curve, 1G black holes of the indicated mass, born with $\chi=0.7$, are spun down
to within $0.18$ of the $|211\rangle$ threshold within $6\,\mathrm{Gyr}$; a window exists for a given
mass where the white region lies above its dashed curve.
The shaded band marks the boson mass preferred by the free-field fit of Ref.~\cite{Kou2026} (90\%).
(a)~Optimistic inputs: GRS~1915+105 keeps $\chi\ge0.95$ for $4.5\times10^{7}\,\mathrm{yr}$, and the
primary of GW241011 is born with $\chi=0.78$ and keeps $\chi\ge0.69$ for $3\times10^{7}\,\mathrm{yr}$.
(b)~Realistic inputs: GRS~1915+105 keeps $\chi\ge0.87$ over its $3\,\mathrm{Gyr}$ age (the $3\sigma$
lower limit $0.92$ minus the $\simeq0.05$ that accretion from its low-mass donor can restore), and
the primary of GW241011 is born with $\chi=0.72$ and has an age of $5\times10^{7}\,\mathrm{yr}$.
All curves use two-level quasi-equilibrium kinetics with relativistic superradiance rates and
leading-order self-interaction rates.}
\label{fig:window}
